% page 412 of the facsimile (image p-411 of the scan)
% block 160.0 x 233.3 mm ; left margin 8.0 mm
\pgc{-12.700mm}{0.000mm}{
\l{\cel{3}404}
\l{}
\l{\sou{ok.}\cel{2}(\sou{arit.)}\cel{3}La cifro \textquotedbl{}8\textquotedbl{}.}
\l{}
\l{\sou{okarino}. (\sou{muziko).}\cel{2}Muzik-instrumento ek terakoto, truiz-}
\l{\cel{3}ita, qua havas la formo di grosa ovo.\cel{2}DEF.}
\l{}
\l{okaziono. I. Cirkonstanco qua eventas oportune. - II. Cir-}
\l{\cel{3}konstanco qua igas rezolvar agor ulo. - EFIRS.}
\l{}
\l{\sou{okro.}\cel{2}Farbo, argilo friabla, flava o reda, (segun lua kom-}
\l{\cel{3}pozanti kemiala). - DEFIRS.}
\l{}
\l{\sou{okta-}.\cel{2}Prefixo ciencala : ok- .}
\l{}
\l{\sou{oktaedro}. (\sou{geom.)}\cel{2}Qua havas ok edri. DEFIS.}
\l{}
\l{\sou{oktanto}. I. (\sou{astr.}) Disto de 45 gradi (okima parto di la}
\l{\cel{3}cirklokurvo) inter du astri. - II. Instrumento ek sekto-}
\l{\cel{3}ro 45-grada, uzata por mezurer la altesi e la disti. DEFIS.}
\l{}
\l{\sou{oktavo}. I. (\sou{muziko)}.\cel{2}Intertempi de ok gradi. - II. (\sou{litur}-}
\l{\cel{3}\sou{gio katol,}) Intertempo de ok dii qua sequas granda festo-}
\l{\cel{3}dio ekleziala e dum qua eventas la memorigo di ta festo.}
\l{\cel{3}- La lasta dio ek ta serio, en qua la oficio esas maxim}
\l{\cel{3}solena.\cel{2}DEFIS.}
\l{}
\l{\sou{oktiliono}. (\sou{arit.)}\cel{2}lO .}
\l{}
\l{\sou{oktobro}.\cel{2}La monato dekesma di la yaro. - DEFIRS.}
\l{}
\l{\sou{oktogono.}\cel{3}(\sou{geom.)}\cel{2}Figuro qua havas ok anguli ed ok lateri.}
\l{\cel{3}- DEFIS.}
\l{}
\l{\sou{oktopodo}. (\sou{zool.)}\cel{3}Sub-klaso de moluski cefalopoda du-bran-}
\l{\cel{3}kia, qua kontenas la pulpi, la argonauti, e c., opoze a}
\l{\cel{3}la subklaso de la dekapodi (kalmari, sepii. e c.) - L.\cel{1}\sou{octo-}}
\l{}
\l{\sou{okulo}.\cel{2}(\sou{anat.})\cel{2}Organo di la vido-fakultato. - DEFIRS.}
\l{}
\l{\sou{okulta}. Di qua la kauzo permanas celata. - DEFIS.}
\l{}
\l{\sou{okultacar}.\cel{2}(\sou{trans.})\cel{2}(\sou{astron.})\cel{2}(\sou{pri astro)}\cel{2}Celar per in-}
\l{\cel{3}terpozar su. - DEFIS.}
\l{}
\l{\sou{okupar}. trans.) I. Posedeskar loko, instalar su en lu, esta-}
\l{\cel{3}blisar su en lu. - II. Posedeskar (ulu) per absorbar lua}
\l{\cel{3}spirito, lua kordio. - III. Konsumar la tempo (di ulu).}
\l{\cel{3}- DEFIS.}
\l{}
\l{\sou{ol(u).}\cel{3}Pronomo personala, triesma persono, singularo, que}
\l{\cel{3}reprezentas irgo altra kam persono.}
\l{}
\l{olda.\cel{2}Qua vivas, od existas, de multa tempo. - E.}
}
